% GENERATED FILE. Edit canonical lessons, then run npm run generate:books.
% canonical-source-hash: fnv1a64:3bd0cbf7d2f0b8ec
% canonical-lessons: SA-C227-paksah, SA-C227-pucchah, SA-C227-srngam, SA-C227-cancuh, SA-C227-grdhrah

\chapter{Wing, Tail, Horn, Beak, Vulture}
\label{ch:sa-wing-tail-horn-beak-vulture}
\hlchaptermodality{227}

\begin{chapteropening}
\textbf{By the end of this chapter:} \emph{I can say a wing, a tail, a horn, a beak and a vulture.}

Complete the last lesson of chapter 227: I can say a wing, a tail, a horn, a beak and a vulture.
\end{chapteropening}

\section[\texorpdfstring{\sk{पक्षः} (pakṣaḥ)}{पक्षः (pakṣaḥ)}]{\sk{पक्षः} (pakṣaḥ) — a wing}

\label{lesson:SA-C227-paksah}



\begin{quote}
Before the new one: say the Sanskrit for a creature, then the Sanskrit for a winged creature, a bird.
\end{quote}

\subsection*{You'll want to know: \sk{पक्षः}}
\textbf{\sk{पक्षः}} — \emph{pakṣaḥ} — \textquotedblleft{}a wing\textquotedblright{}.

\begin{grammarlens}[title={What you've built}]
One more word for this chapter.
\end{grammarlens}

\subsection*{Your turn}
\begin{itemize}
\raggedright
  \item \emph{Say it:} \emph{pakṣaḥ}
  \item \emph{Say it:} \emph{pakṣaḥ}, once more
  \item \emph{Say it:} say \emph{pakṣaḥ}
  \item \emph{Recall:} say the Sanskrit for a creature, then the Sanskrit for a winged creature, a bird, then say \emph{pakṣaḥ} again
\end{itemize}

\begin{tcolorbox}[breakable,colback=teal!4,colframe=teal!35!black,title={Before you move on}]
What does \sk{पक्षः} mean? (\textquotedblleft{}A wing\textquotedblright{}.) Say it once more.
\end{tcolorbox}
\section[\texorpdfstring{\sk{पुच्छः} (pucchaḥ)}{पुच्छः (pucchaḥ)}]{\sk{पुच्छः} (pucchaḥ) — a tail}

\label{lesson:SA-C227-pucchah}



\begin{quote}
Before the new one: say the Sanskrit for a winged creature, a bird, then the Sanskrit for a wing.
\end{quote}

\subsection*{You'll want to know: \sk{पुच्छः}}
\textbf{\sk{पुच्छः}} — \emph{pucchaḥ} — \textquotedblleft{}a tail\textquotedblright{}.

\begin{grammarlens}[title={What you've built}]
One more word for this chapter.
\end{grammarlens}

\subsection*{Your turn}
\begin{itemize}
\raggedright
  \item \emph{Say it:} \emph{pucchaḥ}
  \item \emph{Say it:} \emph{pucchaḥ}, once more
  \item \emph{Say it:} say \emph{pucchaḥ}
  \item \emph{Recall:} say the Sanskrit for a winged creature, a bird, then the Sanskrit for a wing, then say \emph{pucchaḥ} again
\end{itemize}

\begin{tcolorbox}[breakable,colback=teal!4,colframe=teal!35!black,title={Before you move on}]
What does \sk{पुच्छः} mean? (\textquotedblleft{}A tail\textquotedblright{}.) Say it once more.
\end{tcolorbox}
\section[\texorpdfstring{\sk{शृंगम्} (śṛṅgam)}{शृंगम् (śṛṅgam)}]{\sk{शृंगम्} (śṛṅgam) — a horn}

\label{lesson:SA-C227-srngam}



\begin{quote}
Before the new one: say the Sanskrit for a wing, then the Sanskrit for a tail.
\end{quote}

\subsection*{You'll want to know: \sk{शृंगम्}}
\textbf{\sk{शृंगम्}} — \emph{śṛṅgam} — \textquotedblleft{}a horn\textquotedblright{}.

\begin{grammarlens}[title={What you've built}]
One more word for this chapter.
\end{grammarlens}

\subsection*{Your turn}
\begin{itemize}
\raggedright
  \item \emph{Say it:} \emph{śṛṅgam}
  \item \emph{Say it:} \emph{śṛṅgam}, once more
  \item \emph{Say it:} say \emph{śṛṅgam}
  \item \emph{Recall:} say the Sanskrit for a wing, then the Sanskrit for a tail, then say \emph{śṛṅgam} again
\end{itemize}

\begin{tcolorbox}[breakable,colback=teal!4,colframe=teal!35!black,title={Before you move on}]
What does \sk{शृंगम्} mean? (\textquotedblleft{}A horn\textquotedblright{}.) Say it once more.
\end{tcolorbox}
\section[\texorpdfstring{\sk{चञ्चुः} (cañcuḥ)}{चञ्चुः (cañcuḥ)}]{\sk{चञ्चुः} (cañcuḥ) — a beak}

\label{lesson:SA-C227-cancuh}



\begin{quote}
Before the new one: say the Sanskrit for a tail, then the Sanskrit for a horn.
\end{quote}

\subsection*{You'll want to know: \sk{चञ्चुः}}
\textbf{\sk{चञ्चुः}} — \emph{cañcuḥ} — \textquotedblleft{}a beak\textquotedblright{}.

\begin{grammarlens}[title={What you've built}]
One more word for this chapter.
\end{grammarlens}

\subsection*{Your turn}
\begin{itemize}
\raggedright
  \item \emph{Say it:} \emph{cañcuḥ}
  \item \emph{Say it:} \emph{cañcuḥ}, once more
  \item \emph{Say it:} say \emph{cañcuḥ}
  \item \emph{Recall:} say the Sanskrit for a tail, then the Sanskrit for a horn, then say \emph{cañcuḥ} again
\end{itemize}

\begin{tcolorbox}[breakable,colback=teal!4,colframe=teal!35!black,title={Before you move on}]
What does \sk{चञ्चुः} mean? (\textquotedblleft{}A beak\textquotedblright{}.) Say it once more.
\end{tcolorbox}
\section[\texorpdfstring{\sk{गृध्रः} (gṛdhraḥ)}{गृध्रः (gṛdhraḥ)}]{\sk{गृध्रः} (gṛdhraḥ) — a vulture}

\label{lesson:SA-C227-grdhrah}



\begin{quote}
Before the new one: say the Sanskrit for a horn, then the Sanskrit for a beak.
\end{quote}

\subsection*{You'll want to know: \sk{गृध्रः}}
\textbf{\sk{गृध्रः}} — \emph{gṛdhraḥ} — \textquotedblleft{}a vulture\textquotedblright{}.

\begin{grammarlens}[title={What you've built}]
One more word for this chapter.
\end{grammarlens}

\subsection*{Your turn}
\begin{itemize}
\raggedright
  \item \emph{Say it:} \emph{gṛdhraḥ}
  \item \emph{Say it:} \emph{gṛdhraḥ}, once more
  \item \emph{Say it:} say \emph{gṛdhraḥ}
  \item \emph{Recall:} say the Sanskrit for a horn, then the Sanskrit for a beak, then say \emph{gṛdhraḥ} again
\end{itemize}

\begin{tcolorbox}[breakable,colback=teal!4,colframe=teal!35!black,title={Before you move on}]
What does \sk{गृध्रः} mean? (\textquotedblleft{}A vulture\textquotedblright{}.) Say it once more.
\end{tcolorbox}
